\section{Архитектура программного комплекса \texttt{traffic\_core}}

Разработанный программный комплекс представляет собой высокопроизводительный Linux-демон, написанный на стандарте C++26. В отличие от классических монолитных транспортных симуляторов (например, SUMO или MATSim), архитектура \texttt{traffic\_core} базируется на принципе строгой изоляции вычислительных контуров и отказе от сильной связности (Tight Coupling).

\subsection{Функциональная декомпозиция модулей}

Для предотвращения деградации кодовой базы и появления монолитных классов («God Objects»), система физически разделена на изолированные статические библиотеки (модули). Каждый модуль имеет собственный \texttt{CMakeLists.txt} и инкапсулирует свою зону ответственности:

\begin{itemize}
    \item \textbf{\texttt{router}} (Навигационное ядро) — чистый \textit{stateless}-модуль In-Memory вычислений. Отвечает исключительно за алгоритмический обход графа (TD-ALT), поддержку 4-арных куч (Heap) и работу с плоскими массивами формата CSR. Не содержит логики симуляции.
    \item \textbf{\texttt{data\_provider}} (Среда и Симулятор) — модуль, отвечающий за 1D-кинематику агентов (модель жидкости — Fluid Model), продвижение физического времени и трансляцию сырых метрик загруженности дорог.
    \item \textbf{\texttt{decision\_engine}} (Модуль принятия решений, МПР) — интеллектуальный «мозг» системы. Анализирует потоки данных от \texttt{data\_provider}, выявляет зарождающиеся заторы с помощью пространственных индексов и отправляет асинхронные запросы (триггеры) в \texttt{router} для локального или глобального перестроения маршрутов.
    \item \textbf{\texttt{core}} — базовый слой инфраструктуры, предоставляющий пулы потоков (Thread Pool), Lock-Free очереди (на базе \texttt{moodycamel::ConcurrentQueue}) и диспетчеризацию задач.
    \item \textbf{\texttt{daemon}} — точка входа (\texttt{main.cpp}), реализующая паттерн \textit{Composition Root}.
\end{itemize}

\subsection{Inversion of Control и паттерн Composition Root}

Связывание независимых подсистем осуществляется через паттерн внедрения зависимостей (Dependency Injection). Модули общаются исключительно через абстрактные контракты (например, интерфейсы \texttt{IAgentDataProvider} или \texttt{ITelemetryPublisher}), ничего не зная о внутренней реализации друг друга. 

Сборка итогового конвейера симуляции происходит в единственной точке приложения — функции \texttt{main()}, выступающей в роли корня композиции (Composition Root). Оркестратор инициализирует объекты в строгой последовательности от ядра к периферии, формируя следующий цикл (Main Loop):
\begin{enumerate}
    \item \textbf{Шаг симуляции:} \texttt{data\_provider} продвигает физику на $\Delta t$ и формирует массивы изменений (дельт).
    \item \textbf{Шаг телеметрии:} Извлеченные легковесные дельты передаются в сетевой слой ZeroMQ для асинхронной отправки на веб-клиенты.
    \item \textbf{Шаг аналитики:} МПР забирает сгенерированные симулятором события-алерты (например, критическое отклонение \texttt{ETA}).
    \item \textbf{Шаг маршрутизации:} МПР запрашивает у пула потоков альтернативные маршруты и возвращает их агентам.
\end{enumerate}

Данный подход превращает архитектуру в легко расширяемый конструктор, позволяющий в будущем (например, при интеграции машинного обучения) прозрачно подменить математический МПР на нейросетевой инференс без переписывания транспортного ядра.

\subsection{Zero-Cost OOP и Data-Oriented Design}

Классический объектно-ориентированный дизайн (с обильным использованием интерфейсов и виртуальных функций) провоцирует промахи кэша инструкций процессора и блокирует агрессивный инлайнинг (inlining) на этапе компиляции. В горячем цикле маршрутизатора это приводит к недопустимому нарушению бюджета времени в 190 наносекунд на одну итерацию обхода узла.

Для сочетания чистой архитектуры с аппаратной скоростью выполнения применена парадигма абстракций с нулевой стоимостью (Zero-Cost OOP):
\begin{enumerate}
    \item \textbf{Запрет на виртуальные вызовы (vtable):} В методах обхода графа и симуляции физики (\texttt{tick}, \texttt{step}, \texttt{calculate\_path}) полностью исключено динамическое связывание.
    \item \textbf{Инкапсуляция плоских данных:} Тяжелые массивы \texttt{row\_ptr} и \texttt{col\_ind} (формат CSR) инкапсулируются внутри классов как приватные поля \texttt{std::vector}. Внешний интерфейс предоставляет безопасные геттеры, помеченные директивами \texttt{inline} и \texttt{[[nodiscard]]}.
    \item \textbf{Compile-Time полиморфизм:} Компилятор LLVM/GCC полностью вырезает вызов таких методов, заменяя их прямым ассемблерным смещением в ОЗУ. Это гарантирует скорость языка ассемблера при сохранении строгой типизации и безопасности памяти современного C++.
\end{enumerate}

Архитектурная схема взаимодействия модулей представлена в Приложении А (рис. \ref{fig:system_arch}).
